\documentclass[a4paper,12pt]{article}
\usepackage[T2A]{fontenc}
\usepackage[utf8]{inputenc}
\usepackage[russian]{babel}
\usepackage{amsmath,amssymb,mathtools}
\usepackage{PTSans}
\renewcommand{\familydefault}{\sfdefault}
\usepackage{icomma}
\usepackage[a4paper,margin=20mm,top=22mm,bottom=20mm,headheight=26pt,headsep=8mm,footskip=11mm]{geometry}
\usepackage{microtype}
\usepackage{xcolor}
\usepackage{tikz}
\usetikzlibrary{arrows.meta,positioning,shapes.geometric,calc}
\usepackage{tabularx,booktabs,array}
\usepackage{enumitem}
\usepackage{fancyhdr}
\usepackage{setspace}
\usepackage[hidelinks]{hyperref}

\definecolor{accent}{HTML}{247774}
\definecolor{darkaccent}{HTML}{17504F}
\definecolor{bodyink}{HTML}{292F30}
\definecolor{muted}{HTML}{667374}
\definecolor{soft}{HTML}{ECF4F2}
\color{bodyink}
\setstretch{1.19}
\setlength{\parindent}{0pt}
\setlength{\parskip}{0.56em}
\setlist[itemize]{leftmargin=1.45em,topsep=0.1em,itemsep=0.15em}
\setlist[enumerate]{leftmargin=1.85em,topsep=0.25em,itemsep=0.75em}
\newcommand{\key}[1]{\textbf{\color{darkaccent}#1}}
\newcommand{\sectionrule}{\vspace{-0.55em}{\color{accent}\hrule height .6pt}\vspace{0.5em}}
\newcolumntype{C}{>{\centering\arraybackslash}X}
\pagestyle{fancy}
\fancyhf{}
\fancyhead[L]{\small\color{darkaccent}Функция и область определения}
\fancyhead[R]{\small\color{darkaccent}Кирилл Зиновьев}
\fancyfoot[C]{\small\color{muted}\thepage}
\renewcommand{\headrulewidth}{.4pt}
\renewcommand{\headrule}{\hbox to\headwidth{\color{accent}\leaders\hrule height\headrulewidth\hfill}}
\tikzset{
 fstart/.style={ellipse,draw=darkaccent,fill=soft,thick,align=center,minimum width=37mm,minimum height=11mm},
 fact/.style={rectangle,rounded corners=2pt,draw=darkaccent,thick,align=center,minimum height=13mm,text width=53mm,inner sep=4pt},
 fchoice/.style={diamond,aspect=2.1,draw=darkaccent,thick,align=center,text width=38mm,inner sep=2pt},
 farrow/.style={-{Latex[length=2.8mm]},thick,draw=accent}
}

\begin{document}
\thispagestyle{empty}
\begin{center}
\vspace*{24mm}
{\Large\color{darkaccent}\bfseries ЧЕК-ЛИСТ}\\[12mm]
{\huge\bfseries Функция и её значения}\\[5mm]
{\Large Таблицы значений. Область определения}\\[13mm]
{\large Алгебра, 7 класс}\\[3mm]
{\large 17 октября 2026 года}\\[21mm]
{\LARGE\bfseries Кирилл Зиновьев}\\[14mm]
\begin{minipage}{.78\linewidth}\centering
Как подставлять аргумент в формулу, составлять таблицу с заданным шагом и определять, при каких значениях $x$ функция имеет смысл.
\end{minipage}
\vfill
{\large Лёвин Артём Александрович}\\[2mm]
{\normalsize эксклюзивно для Кирилла Зиновьева}\\[13mm]
\begin{tikzpicture}
\draw[accent,line width=1.15pt]
(-6.1,0) .. controls (-4.1,.63) and (-1.9,-.53) .. (0,0)
.. controls (2.2,.66) and (4.2,-.44) .. (6.1,.12);
\end{tikzpicture}
\end{center}
\clearpage

\section*{1. Аргумент и значение функции}
\sectionrule
Если функция задана формулой $y=\dfrac{12}{x}$, буква \key{$x$} обозначает \key{аргумент функции}, а $y$ --- соответствующее \key{значение функции}.

Запись $y(-6)$ читается «значение функции при $x=-6$». Число в скобках указывает, \emph{что подставляем вместо $x$}:
\[
\boxed{y(-6)=\frac{12}{-6}=-2.}
\]
Здесь аргумент равен $-6$, а значение функции равно $-2$. Переход от аргумента к значению один: \key{подставить $x$ в формулу и вычислить}.

\begin{center}
\renewcommand{\arraystretch}{1.35}
\begin{tabularx}{\linewidth}{@{}l*{6}{C}@{}}
\toprule
Аргумент $x$ & $-6$ & $-4$ & $-3$ & $2$ & $6$ & $12$\\
\midrule
Значение $y(x)$ & $-2$ & $-3$ & $-4$ & $6$ & $2$ & $1$\\
\bottomrule
\end{tabularx}
\end{center}

\key{Проверка терминов.} В записи $y(2)=6$ число $2$ --- аргумент, а $6$ --- значение функции.

\section*{2. Подстановка отрицательных чисел}
\sectionrule
Для функции $y=x^2-9$ при $x=-4$ получаем:
\[
y(-4)=\underbrace{(-4)^2}_{16}-9=7.
\]
\key{Скобки сохраняют знак подставляемого числа.} Сравните:
\[
(-4)^2=16,\qquad -4^2=-(4^2)=-16.
\]
Без скобок знак минус \emph{не входит} в основание степени. Когда подставляете отрицательное число в выражение со степенью, берите его в скобки.

\begin{center}
\renewcommand{\arraystretch}{1.3}
\begin{tabularx}{.84\linewidth}{@{}l*{5}{C}@{}}
\toprule
$x$ & $-5$ & $-4$ & $-3$ & $0$ & $2$\\
\midrule
$y=x^2-9$ & $16$ & $7$ & $0$ & $-9$ & $-5$\\
\bottomrule
\end{tabularx}
\end{center}

\section*{3. Таблица значений с заданным шагом}
\sectionrule
Пусть $y=x(x-3{,}5)$, а $x$ изменяется от $0$ до $4$ с шагом $0{,}5$. Значит, начинаем с $0$ и \key{прибавляем $0{,}5$}, пока не дойдём до $4$:
\[
0;\ 0{,}5;\ 1;\ 1{,}5;\ 2;\ 2{,}5;\ 3;\ 3{,}5;\ 4.
\]
Далее вычисляем значения функции для каждого аргумента:
\[
\begin{aligned}
y(0)&=0\cdot(0-3{,}5)=0,\\
y(0{,}5)&=0{,}5\cdot(0{,}5-3{,}5)=-1{,}5.
\end{aligned}
\]
\begin{center}
\renewcommand{\arraystretch}{1.25}
\small
\begin{tabularx}{\linewidth}{@{}l*{9}{C}@{}}
\toprule
$x$ & $0$ & $0{,}5$ & $1$ & $1{,}5$ & $2$ & $2{,}5$ & $3$ & $3{,}5$ & $4$\\
\midrule
$y(x)$ & $0$ & $-1{,}5$ & $-2{,}5$ & $-3$ & $-3$ & $-2{,}5$ & $-1{,}5$ & $0$ & $2$\\
\bottomrule
\end{tabularx}
\end{center}
\key{Не перепутайте:} шаг относится к изменениям \emph{аргумента} $x$. Значения функции вовсе не обязаны изменяться на одинаковую величину.

\section*{4. Область определения функции}
\sectionrule
\key{Область определения} --- все значения аргумента $x$, которые можно подставить в формулу и получить значение функции.

Обозначение $\mathbb R$ означает \key{множество всех действительных чисел}. Для изучаемых формул без ограничений пишем $x\in\mathbb R$.

Если в знаменателе есть выражение с $x$, проверяем, \key{когда оно становится нулём}. Деление на нуль запрещено, поэтому соответствующие аргументы исключают. Это правило относится к изученным видам функций; позднее Вы познакомитесь и с другими ограничениями.

\begin{center}
\renewcommand{\arraystretch}{1.95}
\begin{tabularx}{\linewidth}{@{}p{39mm}X p{50mm}@{}}
\toprule
\textbf{Формула} & \textbf{Проверка знаменателя} & \textbf{Область определения}\\
\midrule
$y=x+5$ & Деления нет. & $x\in\mathbb R$\\
$y=\dfrac{x}{5}$ & Знаменатель $5\ne0$. & $x\in\mathbb R$\\
$y=\dfrac{1}{x-7}$ & $x-7=0\Rightarrow x=7$. & $x\in\mathbb R$, $x\ne7$\\
$y=\dfrac{2}{x+3}$ & $x+3=0\Rightarrow x=-3$. & $x\in\mathbb R$, $x\ne-3$\\
$y=\dfrac{4x-1}{5}$ & Знаменатель $5\ne0$. & $x\in\mathbb R$\\
\bottomrule
\end{tabularx}
\end{center}

\key{Почему $x=7$ исключён?} Для $y=\dfrac1{x-7}$ подстановка семёрки приводит к $\dfrac1{7-7}=\dfrac10$, а делить на нуль нельзя. Все другие действительные значения $x$ допустимы.

\key{Перед ответом проверьте:} правильно ли подставлен аргумент, учтены ли скобки при отрицательном числе, соблюдён ли шаг и исключены ли значения, обнуляющие знаменатель.
\clearpage

\section*{Алгоритм: как найти область определения}
\sectionrule
\vspace{10mm}
\begin{center}
\begin{tikzpicture}[node distance=12mm and 19mm,font=\small]
\node[fstart] (start) {Начало};
\node[fact,below=of start] (formula) {Рассмотрите формулу функции.};
\node[fchoice,below=17mm of formula] (ask) {Есть выражение с $x$ в знаменателе?};
\node[fact,below left=26mm and 13mm of ask,text width=49mm] (no) {Для изученных формул все значения $x$ допустимы: $x\in\mathbb R$.};
\node[fact,below right=26mm and 13mm of ask,text width=51mm] (yes) {Приравняйте знаменатель к нулю и решите уравнение.};
\node[fact,below=16mm of yes,text width=51mm] (exclude) {Исключите найденные значения $x$ из $\mathbb R$.};
\node[fstart,below=25mm of no,minimum width=35mm] (endno) {Ответ записан};
\node[fstart,below=13mm of exclude,minimum width=35mm] (endyes) {Ответ записан};
\draw[farrow] (start)--(formula);
\draw[farrow] (formula)--(ask);
\draw[farrow] (ask.south west) -| node[pos=.25,above left,font=\scriptsize]{нет} (no.north);
\draw[farrow] (ask.south east) -| node[pos=.25,above right,font=\scriptsize]{да} (yes.north);
\draw[farrow] (no)--(endno);
\draw[farrow] (yes)--(exclude);
\draw[farrow] (exclude)--(endyes);
\end{tikzpicture}
\end{center}
\clearpage

\section*{Тренировочный блок. Путь А}
\sectionrule
\textit{10 заданий от простого к более сложному. Записывайте вычисления и конечный ответ; решения задач на этой странице не приведены.}
\begin{enumerate}[label=\textbf{\arabic*.},itemsep=1.05em]
\item Для функции $y=\dfrac{12}{x}$ вычислите $y(-6)$. Назовите аргумент и значение функции.
\item Для $y=\dfrac{12}{x}$ найдите $y(-4)$, $y(3)$ и $y(6)$. Оформите вычисления через запись $y(a)$.
\item Заполните таблицу для $y=\dfrac{12}{x}$:
\begin{center}\renewcommand{\arraystretch}{1.25}
\begin{tabularx}{.91\linewidth}{@{}l*{5}{C}@{}}\toprule
$x$ & $-12$ & $-3$ & $2$ & $4$ & $12$\\\midrule
$y(x)$ & & & & & \\\bottomrule
\end{tabularx}\end{center}
\item Для $y=x^2-9$ вычислите значения функции при $x=-5$; $-4$; $0$; $2$.
\item Ученик получил $y(-3)=-18$ для функции $y=x^2-9$. Найдите ошибку, запишите верное значение и объясните, зачем нужны скобки.
\item Составьте таблицу значений $y=x(x-3)$ для $1\le x\le3$ с шагом $1$.
\item Составьте таблицу значений $y=x(x-3{,}5)$ для $0\le x\le2$ с шагом $0{,}5$.
\item Найдите области определения функций: а) $y=x+5$; б) $y=\dfrac{x}{5}$; в) $y=\dfrac5x$. Объясните различие.
\item Найдите области определения функций: а) $y=\dfrac1{x-7}$; б) $y=\dfrac2{x+3}$; в) $y=\dfrac{4x-1}{5}$.
\item Для $y=\dfrac{12}{x-3}$ найдите область определения и заполните таблицу при $x=-1$; $0$; $2$; $3$; $4$; $5$. Если значение функции не существует, поставьте прочерк и объясните причину.
\end{enumerate}
\clearpage

\section*{Ответы к тренировочному блоку}
\sectionrule
\begin{center}
\renewcommand{\arraystretch}{1.48}
\begin{tabularx}{\linewidth}{@{}p{10mm}X@{}}
\toprule
\textbf{№} & \textbf{Ответ}\\
\midrule
1 & $y(-6)=-2$. Аргумент $-6$, значение функции $-2$.\\
2 & $y(-4)=-3$; $y(3)=4$; $y(6)=2$.\\
3 & Для $x=-12;-3;2;4;12$ соответственно $y=-1;-4;6;3;1$.\\
4 & Для $x=-5;-4;0;2$ соответственно $y=16;7;-9;-5$.\\
5 & $y(-3)=(-3)^2-9=9-9=0$. Ошибка: $-3^2=-9$, но $(-3)^2=9$.\\
6 & $x$: $1;2;3$. Значения $y$: $-2;-2;0$.\\
7 & $x$: $0;0{,}5;1;1{,}5;2$. Значения $y$: $0;-1{,}5;-2{,}5;-3;-3$.\\
8 & а) $x\in\mathbb R$; б) $x\in\mathbb R$; в) $x\in\mathbb R$, $x\ne0$. В пункте в) знаменатель зависит от $x$.\\
9 & а) $x\in\mathbb R$, $x\ne7$; б) $x\in\mathbb R$, $x\ne-3$; в) $x\in\mathbb R$.\\
10 & $x\in\mathbb R$, $x\ne3$. Значения для $x=-1;0;2;3;4;5$ соответственно: $-3;-4;-12;\text{---};12;6$. При $x=3$ деление на нуль.\\
\bottomrule
\end{tabularx}
\end{center}
\vspace{1em}
\key{Как себя проверить.} В заданиях 1--7 повторно подставьте каждый аргумент в формулу. В заданиях 8--10 убедитесь, что ни одно найденное допустимое значение не обращает знаменатель в нуль.
\end{document}
